\chapter{Zbieżność zmiennych losowych i rozkładów}
\label{ch:convergence}

\providecommand{\atlasref}[2]{#2}

\paragraph{Odsyłacze międzytomowe.}
Rodzaje zbieżności wykorzystywane w statystyce
asymptotycznej omówiono w tomie IV, w rozdziale
o asymptotycznej teorii estymacji%
\IfFileExists{probability_volume_iv.aux}{ (rozdział~\ref{IV-ch:asymptotic-estimation})}{}.


Pojęcie granicy występuje w probabilistyce w kilku znaczeniach.
Przy wielokrotnych pomiarach chcielibyśmy wiedzieć, czy
ich średnia zbliży się do pewnej liczby; przy przybliżaniu
rozkładu dwumianowego rozkładem Poissona pytamy natomiast,
czy zbliżają się do siebie \emph{rozkłady}, niekoniecznie
realizacje zmiennych na jednej przestrzeni. Te zadania
wymagają rozróżnienia kilku rodzajów zbieżności.

\section{Zbieżność prawie na pewno, według prawdopodobieństwa i w~\texorpdfstring{$L^p$}{Lp}}

\begin{definicja}
Niech $(\Omega,\mathcal F,\Prob)$ będzie przestrzenią
probabilistyczną, a $X,X_1,X_2,\ldots:\Omega\to\R$
rzeczywistymi zmiennymi losowymi.
\begin{enumerate}
\item Mówimy, że $X_n$ jest zbieżny do $X$ \emph{prawie na pewno}
 i piszemy $X_n\to X$ p.n., gdy
 $\Prob(\{\omega:\lim_{n\to\infty}X_n(\omega)=X(\omega)\})=1$.
\item Mówimy, że $X_n$ jest zbieżny do $X$
 \emph{według prawdopodobieństwa}, $X_n\xrightarrow{\Prob}X$,
 gdy dla każdego $\varepsilon>0$ zachodzi
 $\lim_{n\to\infty}\Prob(|X_n-X|>\varepsilon)=0$.
\item Dla $1\leq p<\infty$ i $X,X_n\in L^p(\Omega,\mathcal F,\Prob)$
 mówimy o \emph{zbieżności w $L^p$}, gdy
 $\E|X_n-X|^p\to0$.
\end{enumerate}
W pierwszym punkcie zbiór, na którym zachodzi zbieżność,
jest mierzalny, ponieważ można go zapisać jako
\[
\bigcap_{r=1}^\infty\bigcup_{m=1}^\infty
\bigcap_{n=m}^\infty\{|X_n-X|\leq1/r\}.
\]
Jeżeli $X_n(\omega)\to X(\omega)$, to dla każdego $r$
nierówność $|X_n(\omega)-X(\omega)|\leq1/r$ zachodzi
od pewnego indeksu. Odwrotnie, dla dowolnego
$\varepsilon>0$ wybieramy $r$ takie, że $1/r<\varepsilon$,
i otrzymujemy od pewnego indeksu odchylenie mniejsze
od $\varepsilon$. Dowodzi to równości zbiorów.
Każdy zbiór po prawej stronie jest mierzalny,
a wszystkie operacje mnogościowe są przeliczalne.
W trzecim punkcie zmienne utożsamiamy w $L^p$ z ich
klasami równoważności względem równości prawie na pewno.
Warunek jest równoważny $\|X_n-X\|_p\to0$, gdzie
$\|Y\|_p=(\E|Y|^p)^{1/p}$, i nie zależy od wyboru
reprezentantów tych klas: dla dwóch wyborów
reprezentantów wszystkich zmiennych wystarczy usunąć
przeliczalną sumę mierzalnych zbiorów zerowych,
na których te wybory się różnią. Poza tą sumą
wartości wszystkich zmiennych i ich różnic są identyczne.
Całki potęg modułów różnic oraz prawdopodobieństwa
zdarzeń odchyleń nie zmieniają się przy modyfikacjach
na mierzalnych zbiorach zerowych. Również warunek
zbieżności punktowej poza zbiorem zerowym pozostaje
niezmieniony.
\end{definicja}

Warunek „prawie na pewno” dotyczy niemal każdej
\emph{trajektorii} $n\mapsto X_n(\omega)$.
Zbieżność według prawdopodobieństwa pozwala na sporadyczne
duże odchylenia, o ile ich prawdopodobieństwo zanika.
Zbieżność w $L^p$ kontroluje natomiast średni koszt
odchylenia podniesionego do potęgi $p$.

Granica według prawdopodobieństwa jest jednoznaczna
z dokładnością do równości prawie na pewno.
Istotnie, jeżeli $X_n\xrightarrow{\Prob}X$ oraz
$X_n\xrightarrow{\Prob}Y$ na tej samej przestrzeni,
to dla $\varepsilon>0$ z nierówności trójkąta mamy
\[
\Prob(|X-Y|>\varepsilon)
\leq\Prob(|X-X_n|>\varepsilon/2)
+\Prob(|X_n-Y|>\varepsilon/2)\longrightarrow0.
\]
Lewa strona nie zależy od $n$, więc jest zerowa.
Ponieważ $\{X\ne Y\}=\bigcup_{r=1}^\infty
\{|X-Y|>1/r\}$, przeliczalna podaddytywność daje
$\Prob(X\ne Y)=0$. Granica prawie na pewno również
jest jednoznaczna w tym sensie: poza sumą dwóch
zbiorów zerowych korzystamy z jednoznaczności
granicy liczbowej.

\begin{twierdzenie}[Podstawowe implikacje]
Niech $X_n,X$ będą rzeczywistymi zmiennymi losowymi
na wspólnej przestrzeni probabilistycznej.
\begin{enumerate}
\item Jeśli $X_n\to X$ prawie na pewno, to
 $X_n\xrightarrow{\Prob}X$.
\item Jeśli $1\leq p<\infty$, $X_n,X\in L^p$ oraz
 $\E|X_n-X|^p\to0$, to $X_n\xrightarrow{\Prob}X$.
\item Jeśli $1\leq q\leq p<\infty$ i $X_n,X\in L^p$ oraz
 $X_n\to X$ w $L^p$, to $X_n\to X$ w $L^q$.
\end{enumerate}
\end{twierdzenie}
\begin{proof}
Dla pierwszej implikacji ustalmy $\varepsilon>0$.
Funkcje charakterystyczne zdarzeń
$\chi_{\{|X_n-X|>\varepsilon\}}$ są mierzalne
i zbiegają do zera na zbiorze prawdopodobieństwa jeden,
na którym $X_n(\omega)\to X(\omega)$.
Są ograniczone przez stałą funkcję $1$, całkowalną,
gdyż $\Prob(\Omega)=1$.
Z \atlasref{ME-08::thm:dct}{twierdzenia o zbieżności zmajoryzowanej}
wynika
\[
\Prob(|X_n-X|>\varepsilon)
=\E\chi_{\{|X_n-X|>\varepsilon\}}\longrightarrow0.
\]
Dla drugiej implikacji do nieujemnej zmiennej
$|X_n-X|^p$ stosujemy nierówność Markowa z progiem
$\varepsilon^p>0$ (rozdział~\ref{ch:expectation}).
Z monotoniczności potęgi na $[0,\infty)$ dostajemy
\[
\Prob(|X_n-X|>\varepsilon)
=\Prob(|X_n-X|^p>\varepsilon^p)
\leq\frac{\E|X_n-X|^p}{\varepsilon^p}\longrightarrow0.
\]
Dla trzeciej implikacji przypadek $q=p$ jest bezpośredni.
Jeżeli $q<p$, to liczby $p/q$ i $p/(p-q)$
są sprzężonymi wykładnikami Höldera. Dla dowolnego
$Y\in L^p$ funkcja $|Y|^q$ należy do $L^{p/q}$,
a stała funkcja $1$ ma normę w $L^{p/(p-q)}$
równą jeden. Z \atlasref{ME-14::thm:holder-14}{nierówności Höldera}
wynika
\[
\E|Y|^q\leq(\E|Y|^p)^{q/p},
\qquad \|Y\|_q\leq\|Y\|_p.
\]
Dowodzi to najpierw inkluzji $L^p\subseteq L^q$,
a następnie, dla $Y=X_n-X$, żądanej zbieżności
$\|X_n-X\|_q\to0$.
\end{proof}

\begin{przyklad}[Zbieżność w prawdopodobieństwie bez zbieżności p.n.]
Na przestrzeni $(\Omega,\mathcal B(\Omega),\Prob)$,
gdzie $\Omega=[0,1)$ i $\Prob=\lambda_1|_{[0,1)}$
jest ograniczeniem miary Lebesgue'a do zbiorów
borelowskich tego przedziału, ustawmy kolejno
wszystkie przedziały diadyczne
$I_{m,j}=[j2^{-m},(j+1)2^{-m})$ dla
$m=0,1,2,\ldots$ oraz $0\leq j<2^m$.
Zdefiniujmy $X_{2^m+j}=\chi_{I_{m,j}}$.
Przedziały są borelowskie, więc wszystkie $X_n$
są zmiennymi losowymi. Każde całkowite $n\geq1$
ma jednoznaczną postać $2^m+j$ dla $0\leq j<2^m$.
Dla $0<\varepsilon<1$ mamy
$\Prob(|X_n|>\varepsilon)=\lambda_1(I_{m,j})=2^{-m}$,
a dla $\varepsilon\geq1$ prawdopodobieństwo wynosi zero.
Ponieważ $m\to\infty$ przy $n\to\infty$, mamy
$X_n\xrightarrow{\Prob}0$. Jednak dla każdego
$\omega\in[0,1)$ i każdego poziomu $m$ dokładnie
jedna z funkcji charakterystycznych przedziałów
tego poziomu wynosi jeden w punkcie $\omega$,
gdyż przedziały tworzą rozłączny podział $[0,1)$.
Dla $m\geq1$ co najmniej jedna wynosi także zero.
Zatem każda trajektoria przyjmuje wartości jeden i zero
nieskończenie wiele razy; $X_n$ nie jest zbieżny do zera
prawie na pewno. Jest to tzw. ciąg piszący na maszynie
(\emph{typewriter sequence}).
Rysunek~\ref{fig:rp1-07-typewriter} przedstawia kolejne poziomy podziału diadycznego.
\end{przyklad}

\begin{figure}[htbp]
\centering
\begin{tikzpicture}[x=7cm,y=.41cm,>=Stealth]
\foreach \n/\m/\j in {1/0/0,2/1/0,3/1/1,4/2/0,5/2/1,6/2/2,7/2/3,8/3/0,9/3/1,10/3/2,11/3/3,12/3/4,13/3/5,14/3/6,15/3/7} {
 \pgfmathsetmacro{\xa}{\j/(2^\m)}
 \pgfmathsetmacro{\xb}{(\j+1)/(2^\m)}
 \fill[blue!55] (\xa,-\n+.12) rectangle (\xb,-\n+.88);
 \draw[gray!35] (0,-\n+.12) rectangle (1,-\n+.88);
 \node[left,font=\scriptsize] at (0,-\n+.5) {$\n$};
}
\draw[->] (0,0)--(1.09,0) node[right] {$\omega$};
\foreach \x/\lab in {0/0,.25/{1/4},.5/{1/2},.75/{3/4},1/1}
 \node[above,font=\scriptsize] at (\x,0) {$\lab$};
\draw[orange!90!black,thick,dashed] (.32,.06)--(.32,-15);
\node[below,font=\small] at (.32,-15) {$\omega=0{,}32$};
\node[left,font=\small] at (-.055,-.5) {$n$};
\end{tikzpicture}
\caption{Pierwsze 15 funkcji ciągu piszącego na maszynie.
Niebieski odcinek w wierszu $n$ oznacza zbiór,
na którym $X_n=1$. Długości odcinków maleją,
ale ustalony punkt $\omega$ należy do dokładnie
jednego przedziału na każdym poziomie diadycznym.}
\label{fig:rp1-07-typewriter}
\end{figure}


\begin{przyklad}[Zbieżność w prawdopodobieństwie bez zbieżności w $L^1$]
Niech $U\sim\Unif(0,1)$ i
$X_n=n\chi_{\{U\leq1/n\}}$ dla $n\geq1$.
Są to zmienne mierzalne i ograniczone (dla każdego
ustalonego $n$), więc należą do $L^p$ dla każdego
$1\leq p<\infty$.
Dla każdego $\varepsilon>0$ i $n>\varepsilon$ mamy
$\{|X_n|>\varepsilon\}=\{U\leq1/n\}$.
Z dystrybuanty rozkładu jednostajnego wynika
$\Prob(|X_n|>\varepsilon)=1/n\to0$.
Zatem $X_n\xrightarrow{\Prob}0$, lecz
$\E|X_n|=n\Prob(U\leq1/n)=1$.
Co więcej, na zdarzeniu $\{U>0\}$ o prawdopodobieństwie
jeden dla każdego $\omega$ od pewnego indeksu
$1/n<U(\omega)$, więc $X_n(\omega)=0$.
Tym samym $X_n\to0$ również prawie na pewno.
Dla $1\leq p<\infty$ mamy
$\E|X_n|^p=n^{p-1}$, co wyklucza zbieżność w $L^p$.
Odchylenia występują coraz rzadziej, ale jednocześnie
stają się coraz większe: średni koszt nie znika.
\end{przyklad}

\begin{twierdzenie}[Kryterium podciągowe zbieżności według prawdopodobieństwa]
Dla rzeczywistych zmiennych losowych $X_n,X$ na wspólnej
przestrzeni probabilistycznej zachodzi
$X_n\xrightarrow{\Prob}X$ wtedy i tylko wtedy, gdy
z każdego podciągu $(X_{n_j})_{j\geq1}$ można wybrać
podciąg $(X_{n_{j_k}})_{k\geq1}$ zbieżny do $X$
prawie na pewno.
\end{twierdzenie}
\begin{proof}
Załóżmy zbieżność według prawdopodobieństwa.
Dla danego podciągu wybieramy rekurencyjnie indeksy
$j_1<j_2<\cdots$ tak, aby
$\Prob(|X_{n_{j_k}}-X|>2^{-k})\leq2^{-k}$.
Jest to możliwe, gdyż dla każdego ustalonego progu
prawdopodobieństwo odchylenia dąży do zera.
Oznaczmy $A_k=\{|X_{n_{j_k}}-X|>2^{-k}\}$.
Mamy $\sum_{k\geq1}\Prob(A_k)\leq\sum_{k\geq1}2^{-k}=1$.
Z pierwszego lematu Borela--Cantellego
(rozdział~\ref{ch:probability}) wynika
$\Prob(\limsup_k A_k)=0$.
Dla każdego $\omega$ spoza tego mierzalnego zbioru zerowego
istnieje $k_0(\omega)$ takie, że dla $k\geq k_0(\omega)$
mamy $|X_{n_{j_k}}(\omega)-X(\omega)|\leq2^{-k}$.
Ponieważ $2^{-k}\to0$, podciąg zbiega do $X$ prawie na pewno.

Odwrotnie, gdyby zbieżność według prawdopodobieństwa
nie zachodziła, istniałyby $\varepsilon,\eta>0$ i
podciąg $(X_{n_j})$ taki, że
$\Prob(|X_{n_j}-X|>\varepsilon)\geq\eta$
dla każdego $j$: jest to negacja warunku zbieżności
nieujemnego ciągu prawdopodobieństw do zera.
Żaden jego podciąg nie mógłby być
zbieżny do $X$ prawie na pewno, bo pierwsza implikacja
poprzedniego twierdzenia wymuszałaby zbieżność
prawdopodobieństw do zera. Otrzymujemy sprzeczność.
\end{proof}

\section{Zbieżność według rozkładu}

\noindent\emph{Powiązania tematyczne w Bibliotece matematycznej:} \atlasref{ME-22::thm:portmanteau-22}{ME-22: twierdzenie Portmanteau}.


Zbieżność według rozkładu nie zakłada nawet istnienia
wspólnej przestrzeni probabilistycznej dla wszystkich
zmiennych. Interesują nas wyłącznie ich miary obrazowe.
Z tego względu jest właściwym językiem dla
przybliżeń Poissona i dla centralnego twierdzenia granicznego.

\begin{definicja}
Niech $X_n$ będą rzeczywistymi zmiennymi losowymi
(niekoniecznie na wspólnej przestrzeni), o dystrybuantach
$F_n$, i niech $X$ będzie zmienną losową o dystrybuancie
$F$. Mówimy, że $X_n$ jest zbieżny do $X$
\emph{według rozkładu}, i piszemy $X_n\xrightarrow{d}X$,
jeśli $F_n(x)\to F(x)$ dla każdego $x\in\R$ będącego
punktem ciągłości funkcji $F$.
Równoważnie mówimy o zbieżności rozkładów
$\mu_n=\Law(X_n)$ do $\mu=\Law(X)$ i piszemy
$\mu_n\Rightarrow\mu$.
Granica jest jednoznaczna jako miara: jeżeli dwie
dystrybuanty $F,G$ spełniają warunek definicji, to
są równe w każdym wspólnym punkcie ciągłości.
Każda z nich ma co najwyżej przeliczalnie wiele
punktów nieciągłości (rozdział~\ref{ch:random-variables}).
Zatem dla dowolnego $x\in\R$ można wybrać
ciąg wspólnych punktów ciągłości $x_k\downarrow x$.
Z prawostronnej ciągłości dostajemy
$F(x)=\lim_k F(x_k)=\lim_k G(x_k)=G(x)$.
Dystrybuanta jednoznacznie wyznacza rozkład.
Nie oznacza to jednoznaczności zmiennej granicznej
jako funkcji na przestrzeni probabilistycznej.
\end{definicja}

Nie wymagamy zbieżności w punktach skoku $F$.
Dla przykładu, zmienna deterministyczna $X_n=1/n$
jest zbieżna do $X=0$ według rozkładu:
dla $x<0$ obie dystrybuanty wynoszą zero,
a dla $x>0$ mamy $F_{X_n}(x)=1$ od pewnego $n$.
Jednak $F_{X_n}(0)=0$ dla wszystkich $n$,
podczas gdy $F_X(0)=1$.
Takie zachowanie nie przeczy definicji.

\begin{twierdzenie}[Zbieżność prawdopodobieństw a zbieżność rozkładów]
Jeśli rzeczywiste zmienne losowe $X_n,X$ są określone
na wspólnej przestrzeni probabilistycznej i
$X_n\xrightarrow{\Prob}X$, to $X_n\xrightarrow{d}X$.
Jeśli $c\in\R$ jest stałą, to zbieżność
$X_n\xrightarrow{d}c$ implikuje
$\Prob(|X_n-c|>\varepsilon)\to0$ dla każdego
$\varepsilon>0$, czyli $X_n\xrightarrow{\Prob}c$.
W tej drugiej części zmienne $X_n$ mogą być określone
na różnych przestrzeniach probabilistycznych
$(\Omega_n,\mathcal F_n,\Prob_n)$, ponieważ $c$
interpretujemy jako funkcję stałą na każdej z nich.
W takim przypadku zapis $X_n\xrightarrow{\Prob}c$
oznacza $\Prob_n(|X_n-c|>\varepsilon)\to0$ dla
każdego $\varepsilon>0$; w dalszym tekście symbol
$\Prob$ pod znakiem granicy oznacza miarę właściwą
zmiennym występującym w danym wyrazie ciągu.
\end{twierdzenie}
\begin{proof}
Niech $x$ będzie punktem ciągłości dystrybuanty $F_X$
i niech $\varepsilon>0$. Oznaczmy $F_n=F_{X_n}$.
Jeżeli $X\leq x-\varepsilon$ oraz $|X_n-X|\leq\varepsilon$,
to $X_n\leq X+\varepsilon\leq x$.
Jeżeli $X_n\leq x$ i $|X_n-X|\leq\varepsilon$,
to $X\leq X_n+\varepsilon\leq x+\varepsilon$.
Z tych dwóch obserwacji otrzymujemy inkluzje zdarzeń
\[
\{X\leq x-\varepsilon\}\cap\{|X_n-X|\leq\varepsilon\}
\subseteq\{X_n\leq x\}
\subseteq\{X\leq x+\varepsilon\}\cup
\{|X_n-X|>\varepsilon\}.
\]
Wynika z nich oszacowanie
\[
F_X(x-\varepsilon)-\Prob(|X_n-X|>\varepsilon)
\leq F_n(x)
\leq F_X(x+\varepsilon)+\Prob(|X_n-X|>\varepsilon).
\]
Istotnie, dla zdarzeń $A,B$ mamy
$\Prob(A\cap B)\geq\Prob(A)-\Prob(B^c)$,
a prawdopodobieństwo sumy zdarzeń nie przekracza
sumy ich prawdopodobieństw.
Ponieważ $X_n\xrightarrow{\Prob}X$, dostajemy
\[
F_X(x-\varepsilon)\leq\liminf_nF_n(x)
\leq\limsup_nF_n(x)\leq F_X(x+\varepsilon).
\]
Po przejściu z $\varepsilon\downarrow0$ i użyciu
ciągłości $F_X$ w $x$ dostajemy $F_n(x)\to F_X(x)$.

Dla drugiego wniosku niech $\varepsilon>0$.
Zdarzenie $\{|X_n-c|>\varepsilon\}$ jest sumą
rozłącznych zdarzeń $\{X_n<c-\varepsilon\}$
oraz $\{X_n>c+\varepsilon\}$, więc
\[
\Prob(|X_n-c|>\varepsilon)
\leq F_n(c-\varepsilon)+1-F_n(c+\varepsilon).
\]
Dystrybuanta $F$ rozkładu skupionego w $c$
jest zerowa dla $x<c$ i równa jeden dla $x\geq c$.
Punkty $c-\varepsilon$ i $c+\varepsilon$ są
punktami jej ciągłości. Oznaczając ponownie
$F_n=F_{X_n}$, z założenia $X_n\xrightarrow{d}c$
otrzymujemy $F_n(c-\varepsilon)\to0$ i
$F_n(c+\varepsilon)\to1$. Otrzymujemy zatem
$\Prob(|X_n-c|>\varepsilon)\to0$.
\end{proof}

\begin{przyklad}[Ten sam rozkład nie oznacza bliskości realizacji]
Niech $X$ przyjmuje wartości $-1$ i $1$ z
prawdopodobieństwami $1/2$, a $X_n=-X$ dla każdego $n$.
Wtedy $\Law(X_n)=\Law(X)$, więc $X_n\xrightarrow{d}X$.
Jednak $|X_n-X|=2$ prawie na pewno, ponieważ
$|X|=1$ prawie na pewno. W szczególności
$\Prob(|X_n-X|>1)=1$ dla każdego $n$, toteż
zbieżność według prawdopodobieństwa do $X$ nie zachodzi.
\end{przyklad}

\begin{lemat}[Całki względem rozkładów zbieżnych]
Jeżeli $X_n\xrightarrow{d}X$, to dla każdej ograniczonej
funkcji ciągłej $h:\R\to\R$ zachodzi
$\E h(X_n)\to\E h(X)$.
W szczególności $\varphi_{X_n}(t)\to\varphi_X(t)$
dla każdego $t\in\R$.
\end{lemat}
\begin{proof}
Oznaczmy $\mu_n=\Law(X_n)$, $\mu=\Law(X)$
oraz $M=\sup_{x\in\R}|h(x)|<\infty$.
Funkcja $h$ jest borelowska, a jej całki względem
wszystkich tych miar są skończone. Jeżeli $M=0$,
teza jest bezpośrednia. Dalej przyjmijmy $M>0$.

Ustalmy $\eta>0$. Z ciągłości miary z dołu dla
rosnących przedziałów $(-R,R]$, których suma jest
równa $\R$, możemy wybrać $R>0$ tak, by
$\mu(\R\setminus(-R,R])<\eta$. Dodatkowo wybieramy
$R$ tak, aby $-R$ i $R$ były punktami ciągłości $F_X$:
wykluczamy co najwyżej przeliczalnie wiele wartości $R$.
Oznaczmy $F_n=F_{X_n}$ i $F=F_X$. Wówczas
\[
\mu_n(\R\setminus(-R,R])=F_n(-R)+1-F_n(R)
\longrightarrow F(-R)+1-F(R)<\eta.
\]
Dla wszystkich dostatecznie dużych $n$ ta wielkość
jest mniejsza niż $2\eta$.

Z jednostajnej ciągłości $h$ na $[-R,R]$ istnieje
$\delta>0$ takie, że $|h(u)-h(v)|<\eta$ dla
$u,v\in[-R,R]$ o $|u-v|<\delta$.
Wybieramy podział $-R=a_0<a_1<\cdots<a_m=R$
o długościach przedziałów mniejszych niż $\delta$,
złożony z punktów ciągłości $F$.
Jest to możliwe, ponieważ zbiór punktów ciągłości
jest gęsty: wystarczy wyjściowy podział o długościach
mniejszych niż $\delta/2$ i odpowiednio małe
przesunięcia jego punktów wewnętrznych.
Definiujemy mierzalną funkcję prostą
\[
s(u)=\sum_{j=1}^m h(a_j)\chi_{(a_{j-1},a_j]}(u).
\]
Na $(-R,R]$ zachodzi $|h-s|<\eta$, a na
dopełnieniu $s=0$, więc $|h-s|\leq M$.
Ponieważ wszystkie $a_j$ są punktami ciągłości $F$,
\[
\int_\R s(u)\,\dd\mu_n(u)
=\sum_{j=1}^m h(a_j)[F_n(a_j)-F_n(a_{j-1})]
\longrightarrow\int_\R s(u)\,\dd\mu(u).
\]
Dla dostatecznie dużych $n$ z nierówności trójkąta
oraz z oszacowań obu ogonów otrzymujemy
\[
\begin{aligned}
\left|\int_\R h\,\dd\mu_n-\int_\R h\,\dd\mu\right|
&\leq\int_\R|h-s|\,\dd\mu_n
+\left|\int_\R s\,\dd\mu_n-\int_\R s\,\dd\mu\right|
+\int_\R|s-h|\,\dd\mu\\
&\leq 2\eta+3M\eta+
\left|\int_\R s\,\dd\mu_n-\int_\R s\,\dd\mu\right|.
\end{aligned}
\]
Przechodząc do granicy górnej po $n$, a następnie
do granicy $\eta\downarrow0$, dostajemy zbieżność
całek. Wzór \eqref{eq:lotus} daje
$\E h(X_n)\to\E h(X)$, niezależnie od przestrzeni,
na których określono zmienne.
Dla $t\in\R$ stosujemy wynik do funkcji
$u\mapsto\cos(tu)$ i $u\mapsto\sin(tu)$.
Obie są ograniczone i ciągłe, stąd z definicji
$\varphi_{X_n}(t)=\E e^{itX_n}$ otrzymujemy
$\varphi_{X_n}(t)\to\varphi_X(t)$.
\end{proof}


\begin{przyklad}[Zbieżność według rozkładu nie oznacza zbieżności według prawdopodobieństwa]
Niech $Z$ przyjmuje wartości $-1$ i $1$ z jednakowymi
prawdopodobieństwami, a $X_n=(-1)^nZ$ i $X=Z$.
Dla każdego $n$ zmienna $X_n$ ma ten sam rozkład co $X$,
a więc $X_n\xrightarrow{d}X$. Jednak dla każdego
nieparzystego $n$ zachodzi $|X_n-X|=2$ prawie na pewno,
więc $\Prob(|X_n-X|>1)=1$.
Wobec tego nie zachodzi $X_n\xrightarrow{\Prob}X$.
Zbieżność według rozkładu nie porównuje wartości
$X_n(\omega)$ i $X(\omega)$ na tych samych trajektoriach;
porównuje wyłącznie miary obrazowe.
\end{przyklad}

\section{Twierdzenie Lévy'ego i przydatne operacje}

Twierdzenie Lévy'ego uzupełnia lemat o całkach względem
rozkładów zbieżnych o implikację odwrotną przy odpowiednim
założeniu o granicy funkcji charakterystycznych.
Jest to jedno z najważniejszych zastosowań tych funkcji:
granicę rozkładów można rozpoznać po granicy ich
transformat Fouriera. Wyjaśnia to, dlaczego funkcje
charakterystyczne będą głównym narzędziem w dowodzie CTG.

Rodzinę borelowskich miar probabilistycznych na $\R$
nazywamy \emph{ciasną}, jeżeli dla każdego $\varepsilon>0$
istnieje $R>0$ takie, że każda miara $\nu$ z tej rodziny
spełnia $\nu(\R\setminus[-R,R])<\varepsilon$.
Równoważnie można wymagać, aby supremum tych wartości
było mniejsze od $\varepsilon$ (w pierwszym warunku
wystarczy użyć $\varepsilon/2$).

\begin{lemat}[Zwartość ciągowa dla ciasnych rodzin na prostej]
Niech $(\mu_n)_{n\geq1}$ będzie ciągiem borelowskich miar
probabilistycznych na $\R$, którego rodzina wyrazów jest ciasna.
Wówczas z każdego jego podciągu można wybrać dalszy
podciąg zbieżny według rozkładu do pewnej borelowskiej
miary probabilistycznej na $\R$.
\end{lemat}
\begin{proof}
\emph{Wybór podciągu.} Niech $F_n$ będą dystrybuantami
$\mu_n$. Uporządkujmy liczby wymierne jako
$\mathbb Q=\{q_1,q_2,\ldots\}$. Ciąg
$(F_n(q_1))_n$ jest ograniczony, zatem ma podciąg
zbieżny. Z tego podciągu wybierzmy następny, na którym
zbiega także $F_n(q_2)$, i kontynuujmy indukcyjnie.
Oznaczmy przez $(n^{(j)}_k)_{k\geq1}$ rosnący ciąg
indeksów wybrany w $j$-tym kroku; każdy jest podciągiem
poprzedniego, a ciąg $(F_{n^{(j)}_k}(q_j))_k$ jest zbieżny.
Przyjmijmy $n_k=n^{(k)}_k$. Wtedy
$n_{k+1}=n^{(k+1)}_{k+1}\geq n^{(k)}_{k+1}>n^{(k)}_k=n_k$.
Dla ustalonego $j$ wszystkie indeksy $n_k$ dla $k\geq j$
należą do ciągu $(n^{(j)}_\ell)_{\ell\geq1}$, zatem
$G(q_j):=\lim_k F_{n_k}(q_j)$ istnieje.
Jeżeli $q,q'\in\mathbb Q$ i $q\leq q'$, to
$F_{n_k}(q)\leq F_{n_k}(q')$ dla wszystkich $k$;
przejście do granicy daje $G(q)\leq G(q')$.
Stąd $G:\mathbb Q\to[0,1]$ jest niemalejąca.

\emph{Konstrukcja dystrybuanty granicznej.}
Dla każdego $x\in\R$ definiujemy
$F(x)=\inf\{G(q):q\in\mathbb Q,\ q>x\}$.
Zbiór wartości, po którym bierzemy infimum, jest
niepusty i zawarty w $[0,1]$. Dla $x<y$ zbiór
wymiernych większych od $y$ zawiera się w zbiorze
wymiernych większych od $x$, więc $F(x)\leq F(y)$.
Dla ciągu $x_j\downarrow x$ mamy $F(x_j)\geq F(x)$.
Z drugiej strony, jeśli ustalimy dowolne wymierne
$q>x$, to dla dostatecznie dużych $j$ zachodzi
$x_j<q$, skąd $F(x_j)\leq G(q)$. Biorąc najpierw
granicę po $j$, a potem infimum po $q>x$,
dostajemy $F(x_j)\to F(x)$. To dowodzi
prawostronnej ciągłości $F$.

\emph{Dlaczego masa nie ucieka do nieskończoności?}
Z ciasności dla ustalonego $\varepsilon>0$
wybieramy $R>0$ tak, aby dla wszystkich $n$
$\mu_n(\R\setminus[-R,R])<\varepsilon$.
Dla każdego wymiernego $q<-R$ mamy
$G(q)\leq\varepsilon$, a dla $q>R$ mamy
$G(q)\geq1-\varepsilon$. Stąd dla $x<-R$
$F(x)\leq\varepsilon$ (wybierając $q\in(x,-R)$),
a dla $x>R$ wszystkie wymierne $q>x$ spełniają
$G(q)\geq1-\varepsilon$, więc również ich infimum
$F(x)\geq1-\varepsilon$.
Ponieważ $\varepsilon$ jest dowolne,
$F(-\infty)=0$ i $F(+\infty)=1$.
Z twierdzenia charakteryzującego dystrybuanty
(rozdział~\ref{ch:random-variables}) wynika istnienie
jedynej borelowskiej miary probabilistycznej $\mu$
o dystrybuancie $F$.

\emph{Zbieżność w punktach ciągłości.}
Niech $x$ będzie punktem ciągłości $F$.
Dla każdego wymiernego $q>x$ monotoniczność
$F_{n_k}$ daje
$\limsup_kF_{n_k}(x)\leq G(q)$, a więc
$\limsup_kF_{n_k}(x)\leq F(x)$.
Z drugiej strony, jeśli $y<x$, wybierzmy
wymierne $q\in(y,x)$. Wtedy
$\liminf_kF_{n_k}(x)\geq G(q)\geq F(y)$,
bo $F(y)$ jest infimum wartości $G$ w punktach
wymiernych większych od $y$. Zatem
$\liminf_kF_{n_k}(x)\geq F(x-)$, gdzie
$F(x-)=\lim_{y\uparrow x}F(y)$.
Ciągłość $F$ w $x$ daje $F(x-)=F(x)$, skąd
$F_{n_k}(x)\to F(x)$. Jest to zbieżność
podciągu według rozkładu. Dowód działa również
dla dowolnego wyjściowego podciągu, gdyż dziedziczy
on ciasność całej rodziny.
\end{proof}

\begin{twierdzenie}[Twierdzenie Lévy'ego o ciągłości]
Niech $(\mu_n)_{n\geq1}$ będzie ciągiem borelowskich miar
probabilistycznych na $\R$, a $\varphi_n$ ich
funkcjami charakterystycznymi. Jeśli dla każdego
$t\in\R$ zachodzi $\varphi_n(t)\to\varphi(t)$,
a funkcja $\varphi$ jest ciągła w zerze, to istnieje
jedyna miara probabilistyczna $\mu$ o funkcji
charakterystycznej $\varphi$ i $\mu_n$ zbiega
według rozkładu do $\mu$.
Odwrotnie, jeżeli $\mu_n\Rightarrow\mu$, to
$\varphi_n(t)\to\varphi_\mu(t)$ dla każdego $t\in\R$,
a $\varphi_\mu$ jest ciągła w zerze.
\end{twierdzenie}
\begin{proof}
Odwrotna implikacja jest treścią lematu
o całkach względem rozkładów zbieżnych, a ciągłość
$\varphi_\mu$ wynika z własności funkcji charakterystycznych
(rozdział~\ref{ch:expectation}).
Dla właściwej części dowodu najpierw wykażemy ciasność.
Dla każdego $n$ realizujemy $X_n$ jako funkcję
identyczności na przestrzeni $(\R,\mathcal B(\R),\mu_n)$,
więc $X_n$ ma rozkład $\mu_n$. Ustalmy $T>0$.
Funkcja $(t,x)\mapsto1-\cos(tx)$ jest ciągła,
nieujemna i ograniczona przez $2$.
Jest mierzalna względem produktowej $\sigma$-algebry
borelowskiej, która pokrywa się z $\sigma$-algebrą
borelowską produktu $[-T,T]\times\R$.
Miara $\lambda_1|_{[-T,T]}\otimes\mu_n$ jest skończona
i ma masę $2T$, więc całka modułu tej funkcji
nie przekracza $4T$. Z
\atlasref{ME-12::thm:fubini-12}{twierdzenia Fubiniego}
i definicji funkcji charakterystycznej otrzymujemy
\[
\begin{aligned}
\frac1{2T}\int_{-T}^T[1-\operatorname{Re}\varphi_n(t)]
\,\dd\lambda_1(t)
&=\int_\R\frac1{2T}\int_{-T}^T(1-\cos(tx))
\,\dd\lambda_1(t)\,\dd\mu_n(x)\\
&=\E\left[1-\frac{\sin(TX_n)}{TX_n}\right].
\end{aligned}
\]
W ostatnim kroku dla $x\ne0$ całka z $\cos(tx)$
wynosi $2\sin(Tx)/x$, a dla $x=0$ wynosi $2T$;
iloraz $\sin(Tx)/(Tx)$ określamy więc w zerze jako $1$.
Jeżeli $|X_n|>2/T$, to
$|\sin(TX_n)/(TX_n)|\leq1/2$.
Ponieważ $|\sin u|\leq |u|$, funkcja
$1-\sin u/u$ z powyższą konwencją jest nieujemna.
Na zdarzeniu $\{|X_n|>2/T\}$ wynosi ona co najmniej
$1/2$. Zatem $\chi_{\{|X_n|>2/T\}}\leq
2(1-\sin(TX_n)/(TX_n))$ i po scałkowaniu mamy
\[
\Prob(|X_n|>2/T)\leq\frac1T
\int_{-T}^T[1-\operatorname{Re}\varphi_n(t)]
\,\dd\lambda_1(t).
\]
Dla ustalonego $T$ funkcje $1-\operatorname{Re}\varphi_n$
są borelowskie, nieujemne i ograniczone przez $2$.
Ich granica punktowa jest borelowska, a stała $2$
jest całkowalna na $[-T,T]$. Z
\atlasref{ME-08::thm:dct}{twierdzenia o zbieżności zmajoryzowanej}
prawa strona zbiega do
$T^{-1}\int_{-T}^T[1-\operatorname{Re}\varphi(t)]\,\dd\lambda_1(t)$.
Ustalmy $\eta>0$. Ponieważ $\varphi(0)=1$
i $\varphi$ jest ciągła w zerze, wybieramy $T>0$ takie,
że $|1-\varphi(t)|<\eta/8$ dla $|t|\leq T$.
Wówczas granica prawej strony jest nie większa niż
$\eta/4$. Istnieje więc $N$ takie, że dla $n\geq N$
prawa strona jest mniejsza niż $\eta/2$.
Dla każdej z miar $\mu_1,\ldots,\mu_{N-1}$ można
wybrać $R_j>0$ tak, by $\mu_j(\R\setminus[-R_j,R_j])<\eta/2$,
gdyż $[-R,R]$ rosną do $\R$.
Biorąc $R$ nie mniejsze niż $2/T$ oraz wszystkie $R_j$,
dostajemy $\sup_n\mu_n(\R\setminus[-R,R])\leq\eta/2<\eta$.
Dowodzi to ciasności.

Z lematu o zwartości wynika, że każdy podciąg $\mu_n$
ma dalszy podciąg zbieżny według rozkładu do pewnej
miary $\nu$. Z lematu o całkach ograniczonych funkcji ciągłych,
zastosowanego do zbieżnego podciągu $(\mu_{n_k})$,
wynika $\varphi_{\mu_{n_k}}(t)\to\varphi_\nu(t)$
dla każdego $t\in\R$. Z założenia cały ciąg
$\varphi_n(t)$ zbiega do $\varphi(t)$, więc
$\varphi_\nu(t)=\varphi(t)$ dla wszystkich $t$.
Istnieje co najmniej jedna taka granica podciągu;
oznaczmy ją przez $\mu$. Ma ona funkcję charakterystyczną
$\varphi$. Z twierdzenia o jednoznaczności rozkładu
wyznaczonego przez funkcję charakterystyczną
(rozdział~\ref{ch:expectation}) wynika, że każda inna
taka granica jest równa $\mu$, a także że $\mu$ jest
jedyną miarą o funkcji charakterystycznej $\varphi$. Gdyby cały ciąg nie zbiegał do
$\mu$, istniałby punkt ciągłości jej dystrybuanty
i podciąg, na którym odchylenie dystrybuanty od granicy
byłoby większe od pewnej dodatniej stałej. Taki podciąg
ma jednak dalszy podciąg zbieżny do $\mu$, co daje
sprzeczność. Zatem $\mu_n\to\mu$ według rozkładu.
\end{proof}

\begin{twierdzenie}[Ciągłe przekształcenia i twierdzenie Slutskiego]
Niech $X_n,X$ będą rzeczywistymi zmiennymi losowymi.
Jeżeli $X_n\xrightarrow{d}X$ i $g:\R\to\R$ jest ciągła,
to $g(X_n)\xrightarrow{d}g(X)$.
Jeżeli $X_n\xrightarrow{d}X$, a dla każdego $n$
rzeczywiste zmienne $X_n,Y_n$ są określone na jednej
przestrzeni probabilistycznej (która może zależeć od $n$)
oraz $\Prob(|Y_n-c|>\varepsilon)\to0$ dla każdego
$\varepsilon>0$ i pewnej stałej
$c\in\R$, to
$X_n+Y_n\xrightarrow{d}X+c$ oraz
$X_nY_n\xrightarrow{d}cX$; jeśli $c\ne0$, także
$X_n/Y_n\xrightarrow{d}X/c$, przy dowolnym określeniu
ilorazu jako skończonej mierzalnej funkcji
na zdarzeniu $\{Y_n=0\}$ (np. przez zero).
Zmienna graniczna $X$ nie musi być określona na tej
samej przestrzeni co pary $(X_n,Y_n)$.
\end{twierdzenie}
\begin{proof}
Korzystamy z charakteryzacji zbieżności według rozkładu
przez całki ograniczonych funkcji ciągłych
(jedną implikację wykazano powyżej, a dowód odwrotnej
podano w sekcji „Dokładna charakteryzacja słabej zbieżności”).
Z ciągłości $g$ wynika jego mierzalność borelowska,
zatem $g(X_n)$ oraz $g(X)$ są zmiennymi losowymi.
Dla ograniczonej ciągłej funkcji $h$ również
$h\circ g$ jest ograniczona i ciągła, więc
$\E h(g(X_n))\to\E h(g(X))$.
Z implikacji odwrotnej charakteryzacji dostajemy
$g(X_n)\xrightarrow{d}g(X)$.

Dla twierdzenia Slutskiego wykażmy najpierw ciasność.
Niech $\eta>0$. Jak w dowodzie lematu o całkach
ograniczonych funkcji ciągłych, wybieramy $R>0$
tak, aby $\pm R$ były punktami ciągłości $F_X$ oraz
$F_X(-R)+1-F_X(R)<\eta/2$.
Wówczas dla dostatecznie dużych $n$
$\Prob(|X_n|>R)\leq F_{X_n}(-R)+1-F_{X_n}(R)<\eta$.
Zwiększenie $R$ pozwala objąć także skończenie wiele
pozostałych rozkładów. Stosując ten argument z $\eta/2$,
możemy zapewnić $\sup_n\Prob(|X_n|>R)<\eta$.

Dla $R,\delta>0$ zdarzenie
$G_n=\{|X_n|\leq R,\ |Y_n-c|\leq\delta\}$
jest mierzalne na przestrzeni wspólnej dla $X_n,Y_n$.
Na nim zachodzą oszacowania
\[
\begin{aligned}
|(X_n+Y_n)-(X_n+c)|&\leq\delta,\\
|X_nY_n-cX_n|&\leq R\delta.
\end{aligned}
\]
Jeżeli ponadto $c\ne0$ i $0<\delta<|c|/2$,
to $|Y_n|\geq |c|-\delta>|c|/2$ na $G_n$, więc
\[
\left|\frac{X_n}{Y_n}-\frac{X_n}{c}\right|
=\frac{|X_n||Y_n-c|}{|Y_n||c|}
\leq\frac{2R\delta}{|c|^2}.
\]
Dla ustalonych $\varepsilon,\eta>0$ wybieramy najpierw
$R$ z ciasności, a następnie $\delta>0$ tak, aby
odpowiednie prawe strony były mniejsze niż $\varepsilon$.
Dla każdej rozważanej różnicy $D_n$ mamy zatem
\[
\Prob(|D_n|>\varepsilon)
\leq\Prob(G_n^c)
\leq\Prob(|X_n|>R)+\Prob(|Y_n-c|>\delta).
\]
Granica górna po $n$ jest nie większa niż $\eta$,
gdyż $Y_n\xrightarrow{\Prob}c$.
Dowolność $\eta$ daje $D_n\xrightarrow{\Prob}0$.
W przypadku ilorazu dowód nie zależy od wybranej
mierzalnej wartości na $\{Y_n=0\}$, ponieważ ten
zbiór nie przecina $G_n$.

Pozostaje uzasadnić, że takie małe zaburzenia
nie zmieniają granicy rozkładów. Niech dla każdego $n$
zmienne rzeczywiste $A_n,B_n$ będą określone na tej
samej przestrzeni, niech
$\Prob(|A_n-B_n|>\varepsilon)\to0$ dla każdego
$\varepsilon>0$ oraz $B_n\xrightarrow{d}B$.
Stosując wcześniejsze oszacowanie dystrybuant
z $A_n$ w roli $X_n$ i $B_n$ w roli $X$, dostajemy
\[
F_{B_n}(x-\varepsilon)-\Prob(|A_n-B_n|>\varepsilon)
\leq F_{A_n}(x)
\leq F_{B_n}(x+\varepsilon)+\Prob(|A_n-B_n|>\varepsilon).
\]
Jeśli $x$ jest punktem ciągłości dystrybuanty $B$,
wybieramy $\varepsilon$ tak, aby także
$x-\varepsilon$ i $x+\varepsilon$ były punktami
ciągłości (dystrybuanta ma co najwyżej przeliczalnie
wiele skoków). Można wybrać ciąg takich $\varepsilon>0$ malejący do zera,
gdyż wykluczamy co najwyżej przeliczalnie wiele liczb.
Dla każdego z tych progów przechodzimy najpierw
do granicy dolnej i górnej po $n$, uzyskując
\[
F_B(x-\varepsilon)\leq\liminf_n F_{A_n}(x)
\leq\limsup_n F_{A_n}(x)\leq F_B(x+\varepsilon).
\]
Ciągłość $F_B$ w $x$ przy $\varepsilon\downarrow0$
daje $F_{A_n}(x)\to F_B(x)$.
Z pierwszej części twierdzenia dla ciągłych funkcji
$x\mapsto x+c$, $x\mapsto cx$ i (gdy $c\ne0$)
$x\mapsto x/c$ wynika zbieżność odpowiednich
przekształceń $X_n$ według rozkładu.
Stosując właśnie dowiedzioną własność małych zaburzeń
do każdej z trzech różnic, otrzymujemy tezę.
\end{proof}

\section{Jednostajna całkowalność a zbieżność wartości oczekiwanych}

\noindent\emph{Powiązania tematyczne w Bibliotece matematycznej:} \atlasref{ME-08::thm:vitali-l1}{ME-08: twierdzenie Vitaliego o zbieżności w normie}.


Zbieżność rozkładów i zbieżność według prawdopodobieństwa
nie muszą powodować zbieżności wartości oczekiwanych.
Wcześniejszy przykład $X_n=n\chi_{\{U\leq1/n\}}$,
gdzie $U\sim\Unif(0,1)$, pokazuje mechanizm:
wkład do pierwszego momentu pochodzi z coraz rzadszych,
ale coraz większych wartości. Dla każdego $K>0$
i całkowitego $n>K$ mamy
$\E[|X_n|\chi_{\{|X_n|>K\}}]=1$.
Zatem jednostajne ograniczenie pierwszych momentów
(nawet $\E|X_n|=1$ dla wszystkich $n$) nie zapewnia
zanikania ogonów tych momentów. Odpowiednim warunkiem wykluczającym
taką sytuację jest jednostajna całkowalność.

\begin{definicja}
Niepustą rodzinę całkowalnych rzeczywistych zmiennych
losowych $\mathcal H$ (niekoniecznie określonych na tej
samej przestrzeni probabilistycznej) nazywamy
\emph{jednostajnie całkowalną}, gdy
\[
\lim_{K\to\infty}\sup_{Y\in\mathcal H}
\E[|Y|\chi_{\{|Y|>K\}}]=0.
\]
Każdą wartość oczekiwaną obliczamy względem miary
probabilistycznej właściwej danej zmiennej. Ze wzoru
\eqref{eq:lotus} dla funkcji nieujemnej wynika
\[
\E[|Y|\chi_{\{|Y|>K\}}]
=\int_\R |y|\chi_{\{|y|>K\}}(y)\,\dd\mu_Y(y),
\qquad \mu_Y=\Law(Y).
\]
Warunek zależy zatem wyłącznie od rodziny rozkładów.
Supremum jest nierosnącą funkcją $K>0$, więc granica
istnieje w $[0,\infty]$.
\end{definicja}

\begin{stwierdzenie}[Kryterium momentowe]
Jeżeli dla pewnego $p>1$ rodzina $\mathcal H$ spełnia
$\sup_{Y\in\mathcal H}\E|Y|^p<\infty$,
to jest jednostajnie całkowalna.
\end{stwierdzenie}
\begin{proof}
Oznaczmy $C=\sup_{Y\in\mathcal H}\E|Y|^p<\infty$.
Każda zmienna jest całkowalna, ponieważ
$|Y|\leq1+|Y|^p$ i miara przestrzeni wynosi jeden.
Dla $K>0$ na zdarzeniu $\{|Y|>K\}$ mamy
$|Y|^{p-1}>K^{p-1}$. Stąd na całej przestrzeni
$|Y|\chi_{\{|Y|>K\}}\leq K^{1-p}|Y|^p$.
Całkując i biorąc supremum, otrzymujemy
\[
0\leq\sup_{Y\in\mathcal H}\E[|Y|\chi_{\{|Y|>K\}}]
\leq CK^{1-p}\longrightarrow0,
\]
gdyż $1-p<0$.
\end{proof}

\begin{twierdzenie}[Zbieżność Vitaliego w $L^1$]
Niech $X_n,X$ będą rzeczywistymi zmiennymi losowymi
na jednej przestrzeni probabilistycznej.
Jeśli $X_n\xrightarrow{\Prob}X$ i rodzina
$\{X_n:n\geq1\}$ jest jednostajnie całkowalna,
to $X\in L^1$ oraz $\E|X_n-X|\to0$.
W szczególności $\E X_n\to\E X$.
\end{twierdzenie}
\begin{proof}
Dla $K>0$ wprowadzamy funkcję obcięcia
$\tau_K(x)=\max(-K,\min(x,K))$.
Jest ona ciągła, ograniczona przez $K$ co do modułu
i spełnia $|\tau_K(x)-\tau_K(y)|\leq|x-y|$,
ponieważ każde z odwzorowań $x\mapsto\min(x,K)$
i $x\mapsto\max(x,-K)$ ma stałą Lipschitza $1$.
Oznaczmy
$u(K)=\sup_n\E[|X_n|\chi_{\{|X_n|>K\}}]$.
Z jednostajnej całkowalności istnieje $K_0>0$
takie, że $u(K_0)\leq1$. Rozdzielając całkę, mamy
\[
\E|X_n|
=\E[|X_n|\chi_{\{|X_n|\leq K_0\}}]
 +\E[|X_n|\chi_{\{|X_n|>K_0\}}]
\leq K_0+u(K_0)\leq K_0+1.
\]
W szczególności $M:=\sup_n\E|X_n|<\infty$.
Z kryterium podciągowego, zastosowanego do całego
ciągu, wybieramy podciąg $X_{n_j}\to X$ prawie na pewno.
Z ciągłości modułu mamy $|X_{n_j}|\to|X|$ prawie na pewno.
Wszystkie te funkcje są mierzalne i nieujemne, więc
\atlasref{ME-08::thm:fatou}{lemat Fatou} daje
\[
\E|X|\leq\liminf_j\E|X_{n_j}|\leq M<\infty.
\]
Zatem $X\in L^1$.
Dla każdego $x\in\R$ zachodzi
$|x-\tau_K(x)|=\max(|x|-K,0)\leq
|x|\chi_{\{|x|>K\}}(x)$.
Zatem $\sup_n\E|X_n-\tau_K(X_n)|\leq u(K)\to0$.
Dla ustalonego $K>0$ niech
$W_n=|\tau_K(X_n)-\tau_K(X)|$.
Mamy $0\leq W_n\leq2K$ oraz $W_n\leq|X_n-X|$
z własności Lipschitza, więc $W_n\xrightarrow{\Prob}0$.
Dla dowolnego $\varepsilon>0$ rozdzielamy całkę:
\[
\E W_n=\E[W_n\chi_{\{W_n\leq\varepsilon\}}]
+\E[W_n\chi_{\{W_n>\varepsilon\}}]
\leq\varepsilon+2K\Prob(W_n>\varepsilon).
\]
Stąd $\limsup_n\E W_n\leq\varepsilon$ dla każdego
$\varepsilon>0$, a więc $\E W_n\to0$.
Ponadto $|X-\tau_K(X)|\leq |X|$ oraz
$|X-\tau_K(X)|\to0$ punktowo przy $K\to\infty$.
Z twierdzenia o zbieżności zmajoryzowanej zastosowanego
do całkowitych $K\to\infty$ otrzymujemy
$\E|X-\tau_K(X)|\to0$; dla dowolnych rzeczywistych
$K\to\infty$ wniosek wynika z monotoniczności
$K\mapsto\max(|X|-K,0)$.
Z nierówności trójkąta otrzymujemy
\[
\E|X_n-X|\leq u(K)+\E W_n+\E|X-\tau_K(X)|.
\]
Dla ustalonego $K$ przechodzimy do granicy górnej
po $n$, co daje
$\limsup_n\E|X_n-X|\leq u(K)+\E|X-\tau_K(X)|$.
Następnie przechodzimy z $K\to\infty$; prawa strona
dąży do zera. Otrzymujemy zbieżność w $L^1$.
Wreszcie z liniowości wartości oczekiwanej i nierówności
$|\E X_n-\E X|\leq\E|X_n-X|$ wynika zbieżność średnich.
\end{proof}

\begin{uwaga}
Implikacje między rodzajami zbieżności mają postać
$L^p\Rightarrow\Prob$ dla $1\leq p<\infty$
i p.n.$\Rightarrow\Prob\Rightarrow d$. Brak strzałki między zbieżnością
prawie na pewno a zbieżnością w $L^p$ bez dodatkowych
założeń. Zbieżność w rozkładzie do \emph{stałej}
implikuje jednak zbieżność według prawdopodobieństwa
do tej stałej.
\end{uwaga}

% Rozszerzenia audytowe rozdzialow 7--9.

\begin{przyklad}[Ta sama granica rozkładów, różne granice średnich]
Niech $U\sim\Unif(0,1)$ i
$X_n=n\chi_{\{U<1/n\}}$. Wtedy na zdarzeniu
$\{U>0\}$ o prawdopodobieństwie jeden dla każdego
$\omega$ od pewnego indeksu $1/n<U(\omega)$.
Zatem $X_n\to0$ prawie na pewno, a z poprzednich
twierdzeń również według rozkładu.
Mimo to $\E X_n=n\Prob(U<1/n)=1$ dla wszystkich $n$,
ponieważ rozkład jednostajny nie ma mas punktowych.
Różnica nie przeczy zbieżności słabej: funkcja
$x\mapsto x$ jest nieograniczona, więc nie jest
jedną z funkcji testowych w charakteryzacji przez
całki z funkcji ograniczonych i ciągłych.
Jednostajna całkowalność zapobiega dokładnie temu
rodzajowi ucieczki masy pierwszych momentów.
\end{przyklad}

\section{Kryteria zbieżności prawie na pewno}

\begin{twierdzenie}[Kryterium sumowalnych odchyleń]
Niech $X,X_1,X_2,\ldots$ będą rzeczywistymi zmiennymi
losowymi na jednej przestrzeni probabilistycznej.
Załóżmy, że dla każdego
$\varepsilon>0$ zachodzi
$\sum_{n=1}^\infty\Prob(|X_n-X|>\varepsilon)<\infty$.
Wówczas $X_n\to X$ prawie na pewno.
\end{twierdzenie}
\begin{proof}
Dla każdej liczby całkowitej $r\geq1$ połóżmy
$A_{n,r}=\{|X_n-X|>1/r\}$. Z mierzalności zmiennych
wynika mierzalność tych zdarzeń, a z założenia
$\sum_n\Prob(A_{n,r})<\infty$.
Z pierwszego lematu Borela--Cantellego
(rozdział~\ref{ch:probability}) otrzymujemy
$\Prob(\limsup_n A_{n,r})=0$.
Zbiór
\[
N=\bigcup_{r=1}^\infty\limsup_{n\to\infty} A_{n,r}
\]
jest mierzalny i ma prawdopodobieństwo zero
z przeliczalnej podaddytywności miary.
Jeżeli $\omega\in\Omega\setminus N$, to dla każdego
$r\geq1$ istnieje $m=m(\omega,r)$ takie, że
$\omega\notin A_{n,r}$ dla wszystkich $n\geq m$,
czyli $|X_n(\omega)-X(\omega)|\leq1/r$.
Dla dowolnego $\varepsilon>0$ wybieramy $r$ tak,
aby $1/r<\varepsilon$; wtedy dla $n\geq m(\omega,r)$
mamy $|X_n(\omega)-X(\omega)|<\varepsilon$.
Zatem $X_n(\omega)\to X(\omega)$ dla wszystkich
$\omega\in\Omega\setminus N$, co dowodzi tezy.
Wystarczyło w istocie założyć sumowalność prawdopodobieństw
dla przeliczalnej rodziny progów $1/r$; niezależność
zdarzeń nie jest potrzebna.
\end{proof}

\begin{uwaga}
Z kryterium wynika również, że
$\sum_{n=1}^\infty\E|X_n-X|^p<\infty$ dla pewnego
rzeczywistego $p>0$ implikuje zbieżność prawie na pewno.
Istotnie, dla każdego $\varepsilon>0$ nierówność Markowa
z progiem $\varepsilon^p$ daje
\[
\sum_{n=1}^\infty\Prob(|X_n-X|>\varepsilon)
\leq\varepsilon^{-p}\sum_{n=1}^\infty\E|X_n-X|^p
<\infty.
\]
Nie trzeba zakładać $p\geq1$, gdyż stosujemy nierówność
do nieujemnej zmiennej $|X_n-X|^p$, a nie korzystamy
z własności normy w $L^p$.

Odwrotność kryterium sumowalnych odchyleń nie zachodzi.
Niech $U\sim\Unif(0,1)$, $Y_n=\chi_{\{U\leq1/n\}}$
i $Y=0$. Na zdarzeniu $\{U>0\}$ o prawdopodobieństwie
jeden każda trajektoria $Y_n$ jest od pewnego miejsca
zerowa. Mimo to dla $0<\varepsilon<1$ mamy
$\Prob(|Y_n|>\varepsilon)=1/n$, a szereg harmoniczny
jest rozbieżny. Również dla każdego $p>0$ zachodzi
$\E|Y_n|^p=1/n$, więc warunek sumowalności momentów
nie jest konieczny do zbieżności prawie na pewno.
W tym przykładzie zdarzenia odchyleń są malejące
i nie można zastosować drugiego lematu Borela--Cantellego
bez sprawdzenia niezależności.
\end{uwaga}

\begin{przyklad}[Zbieżność p.n. bez zbieżności w $L^1$]
Wróćmy do zmiennych $X_n=n\chi_{\{U\leq1/n\}}$,
gdzie $U\sim\Unif(0,1)$.
Na zdarzeniu $\{U>0\}$ każda trajektoria jest
od pewnego miejsca zerowa, więc $X_n\to0$ p.n.
Jednak dla $1\leq p<\infty$
\[
\|X_n\|_p=(n^p\Prob(U\leq1/n))^{1/p}
=n^{1-1/p}.
\]
W szczególności norma w $L^1$ jest stale równa jeden,
a dla $p>1$ normy rosną do nieskończoności.
Zbieżność prawie na pewno nie daje zatem zbieżności
w żadnym z tych $L^p$, nawet gdy każdy wyraz i granica
należą do wszystkich rozważanych przestrzeni.
\end{przyklad}

\begin{przyklad}[Zbieżność w każdym $L^p$, lecz nie p.n.]
Dla ciągu funkcji charakterystycznych przedziałów
diadycznych z przykładu z \emph{ciągiem piszącym na maszynie}
mamy $|X_{2^m+j}|^p=X_{2^m+j}$ dla każdego $p>0$.
Dla każdego ustalonego $1\leq p<\infty$ otrzymujemy
\[
\|X_{2^m+j}\|_p
=(\lambda_1(I_{m,j}))^{1/p}=2^{-m/p}\longrightarrow0
\qquad (m\to\infty),
\]
jednostajnie względem $0\leq j<2^m$.
Ponieważ każdy indeks $n\geq1$ ma tę postać,
$X_n\to0$ w $L^p$ dla każdego takiego $p$.
Jednocześnie, jak wcześniej wykazano, każda trajektoria
przyjmuje wartości zero i jeden nieskończenie wiele razy.
Zatem nawet zbieżność we wszystkich $L^p$ o skończonym
wykładniku nie implikuje zbieżności prawie na pewno.
\end{przyklad}

\begin{przyklad}[Zbieżność według rozkładu bez zbieżności w $L^1$ mimo jednostajnej całkowalności]
Niech $U\sim\Unif(0,1)$ i $X_n=1-U$, $X=U$.
Dla $0<x<1$ mamy
\[
\Prob(X_n\leq x)=\Prob(U\geq1-x)=x,
\]
ponieważ rozkład jednostajny nie ma mas punktowych.
Dla $x\leq0$ dystrybuanta $X_n$ jest zerowa,
a dla $x\geq1$ równa jeden. Zatem
$\Law(X_n)=\Law(X)$ i $X_n\xrightarrow{d}X$.
Jednak
\[
\begin{aligned}
\Prob(|X_n-X|>1/4)
&=\Prob(U<3/8)+\Prob(U>5/8)\\
&=3/8+3/8=3/4.
\end{aligned}
\]
Zbieżność według prawdopodobieństwa nie zachodzi.
Ponadto $|X_n|\leq1$ prawie na pewno, więc rodzina
$\{X_n:n\geq1\}$ jest jednostajnie całkowalna
(jej całki ogonowe są zerowe dla $K\geq1$).
Z wzoru \eqref{eq:lotus} otrzymujemy
\[
\E|X_n-X|
=\int_0^1|1-2u|\,\dd\lambda_1(u)
=\int_0^{1/2}(1-2u)\,\dd\lambda_1(u)
+\int_{1/2}^1(2u-1)\,\dd\lambda_1(u)=\frac12.
\]
Przykład pokazuje, dlaczego w twierdzeniu Vitaliego
nie można zastąpić zbieżności według prawdopodobieństwa
zbieżnością według rozkładu, nawet przy ograniczonych
zmiennych na wspólnej przestrzeni.
\end{przyklad}

\section{Dokładna charakteryzacja słabej zbieżności}

\begin{twierdzenie}[Kryterium funkcji ograniczonych ciągłych]
Dla ciągu borelowskich miar probabilistycznych $(\mu_n)$
na $\R$ oraz borelowskiej miary probabilistycznej $\mu$
na $\R$, o dystrybuantach odpowiednio $F_n$ i $F$,
następujące warunki są równoważne:
\begin{enumerate}
\item dystrybuanty $F_n$ zbiegają do $F$ w każdym
punkcie ciągłości $F$;
\item dla każdej ograniczonej ciągłej funkcji
$h:\R\to\R$ mamy
$\int_\R h(u)\,\dd\mu_n(u)\to\int_\R h(u)\,\dd\mu(u)$.
\end{enumerate}
\end{twierdzenie}
\begin{proof}
Implikacja (1)$\Rightarrow$(2) została wykazana w lemacie
o całkach względem rozkładów zbieżnych: każdą miarę $\nu$
można zrealizować jako rozkład identyczności na przestrzeni
$(\R,\mathcal B(\R),\nu)$, a dowód lematu dotyczy
bezpośrednio całek względem tych rozkładów.
Dla kierunku odwrotnego ustalmy punkt ciągłości $x$
dystrybuanty $F$ oraz $\delta>0$. Definiujemy funkcje
\[
h^-_\delta(u)=
\begin{cases}
1,&u\leq x-\delta,\\
(x-u)/\delta,&x-\delta<u<x,\\
0,&u\geq x,
\end{cases}
\qquad
h^+_\delta(u)=
\begin{cases}
1,&u\leq x,\\
(x+\delta-u)/\delta,&x<u<x+\delta,\\
0,&u\geq x+\delta.
\end{cases}
\]
Obie funkcje są ciągłe (w punktach sklejenia wartości
są zgodne z granicami jednostronnymi) i przyjmują
wartości w $[0,1]$. Punktowo mamy
\[
\chi_{(-\infty,x-\delta]}\leq h^-_\delta
\leq\chi_{(-\infty,x]}\leq h^+_\delta
\leq\chi_{(-\infty,x+\delta]}.
\]
Całkując środkowe nierówności względem $\mu_n$, otrzymujemy
\[
\int_\R h^-_\delta(u)\,\dd\mu_n(u)
\leq F_n(x)\leq
\int_\R h^+_\delta(u)\,\dd\mu_n(u).
\]
Z warunku (2) obie całki graniczne istnieją i są
odpowiednimi całkami względem $\mu$. Całkując zewnętrzne
nierówności względem $\mu$, uzyskujemy
\[
\begin{aligned}
F(x-\delta)&\leq\int_\R h^-_\delta(u)\,\dd\mu(u)
\leq\liminf_nF_n(x)\\
&\leq\limsup_nF_n(x)
\leq\int_\R h^+_\delta(u)\,\dd\mu(u)
\leq F(x+\delta).
\end{aligned}
\]
Gdy $\delta\downarrow0$, lewa i prawa strona
dążą do $F(x)$ na mocy ciągłości w $x$.
Otrzymujemy $F_n(x)\to F(x)$ w każdym punkcie
ciągłości $F$, czyli $\mu_n\Rightarrow\mu$.
\end{proof}
Równoważność pozwala utożsamić zbieżność według rozkładu
na prostej ze słabą zbieżnością miar probabilistycznych
w sensie całek ograniczonych funkcji ciągłych
(por. \atlasref{ME-22::thm:portmanteau-22}{twierdzenie Portmanteau}).

\begin{przyklad}[Ograniczenie funkcji testowej ma znaczenie]
Niech $U\sim\Unif(0,1)$, $X_n=n\chi_{\{U\leq1/n\}}$
i $X=0$. Jak wykazano wcześniej, $X_n\to0$ prawie
na pewno, a zatem również według rozkładu.
Dla każdej ograniczonej ciągłej funkcji $h$ możemy
także sprawdzić wprost, że
\[
\E h(X_n)=\frac1n h(n)+\left(1-\frac1n\right)h(0)
\longrightarrow h(0),
\]
ponieważ $|h(n)|\leq\sup_{x\in\R}|h(x)|<\infty$.
Dla nieograniczonej funkcji ciągłej $h(x)=x$
ten sam wzór daje jednak $\E h(X_n)=1$,
podczas gdy $\E h(X)=0$.
Nie można więc usunąć ograniczoności testów
bez dodatkowej kontroli ogonów.
\end{przyklad}

\section{Zbieżność wartości oczekiwanych przy zbieżności rozkładów}

\begin{twierdzenie}[Jednostajna całkowalność i zbieżność słaba]
Niech $X_n\xrightarrow{d}X$, gdzie zmienne mogą być
określone na różnych przestrzeniach. Jeśli rodzina
$\{X_n:n\geq1\}$ jest jednostajnie całkowalna
w sensie wcześniejszej definicji (zależnej tylko od rozkładów),
to $X\in L^1$ i $\E X_n\to\E X$.
\end{twierdzenie}
\begin{proof}
Oznaczmy $u(K)=\sup_n\E[|X_n|\chi_{\{|X_n|>K\}}]$.
Z jednostajnej całkowalności mamy $u(K)\to0$.
Jak w dowodzie twierdzenia Vitaliego, dla $K_0>0$
takiego, że $u(K_0)\leq1$, otrzymujemy
$\E|X_n|\leq K_0+u(K_0)\leq K_0+1$.
Zatem $M:=\sup_n\E|X_n|<\infty$.
Dla każdego $K>0$ funkcja $x\mapsto\min(|x|,K)$
jest ograniczona i ciągła. Z kryterium funkcji testowych
oraz wzoru \eqref{eq:lotus} wynika
\[
\E\min(|X|,K)=\lim_n\E\min(|X_n|,K)\leq M.
\]
Dla całkowitych $K\to\infty$ nieujemne zmienne
$\min(|X|,K)$ rosną punktowo do $|X|$.
Z \atlasref{ME-08::thm:mct}{twierdzenia o zbieżności monotonicznej}
otrzymujemy $\E|X|\leq M<\infty$, czyli $X\in L^1$.

Niech $\tau_K(x)=\max(-K,\min(x,K))$ będzie funkcją
obcięcia używaną w dowodzie twierdzenia Vitaliego.
Jest ona ograniczona i ciągła, a zatem
$\E\tau_K(X_n)\to\E\tau_K(X)$.
Punktowo mamy
$|x-\tau_K(x)|\leq|x|\chi_{\{|x|>K\}}(x)$.
Ponieważ wszystkie rozważane zmienne są całkowalne,
z liniowości wartości oczekiwanej na każdej z ich
przestrzeni i nierówności trójkąta otrzymujemy
\[
|\E X_n-\E X|\leq
\E|X_n-\tau_K(X_n)|
+|\E\tau_K(X_n)-\E\tau_K(X)|
+\E|X-\tau_K(X)|.
\]
Pierwszy składnik jest nie większy niż $u(K)$,
a trzeci dąży do zera przy całkowitych $K\to\infty$:
funkcje $|X-\tau_K(X)|$ zbiegają punktowo do zera
i są zmajoryzowane przez całkowalną funkcję $|X|$,
więc korzystamy z \atlasref{ME-08::thm:dct}{twierdzenia o zbieżności zmajoryzowanej}.
Dla każdego ustalonego $K$ środkowy składnik dąży do zera.
Przechodząc do granicy górnej po $n$, dostajemy
\[
0\leq\limsup_n|\E X_n-\E X|
\leq u(K)+\E|X-\tau_K(X)|.
\]
Przy $K\to\infty$ po liczbach całkowitych prawa
strona dąży do zera, co kończy dowód.
\end{proof}

\begin{uwaga}
Twierdzenie nie zapewnia zbieżności w $L^1$.
Jeżeli zmienne są określone na różnych przestrzeniach,
różnica $X_n-X$ nie jest w ogóle określona bez wyboru
wspólnej realizacji. Nawet przy wspólnej przestrzeni
wniosek o zbieżności w $L^1$ nie wynika z samych założeń
tego twierdzenia: w poprzednim przykładzie $X_n=1-U$,
$X=U$ mamy jednostajną całkowalność, zbieżność według
rozkładu i $\E X_n=\E X=1/2$, ale
$\E|X_n-X|=1/2$.
Jeżeli zamiast zbieżności według rozkładu założymy
zbieżność według prawdopodobieństwa na wspólnej
przestrzeni, to twierdzenie Vitaliego daje zbieżność
w $L^1$.
\end{uwaga}

% Koniec dodatkow do rozdzialu 7.

